\subsection{Ochrana dat}
\label{subsec:llm-data-security}

Vstupní data odesílaná do LLM mohou obsahovat citlivé informace jakou jsou studentské zdrojové kódy, interní zadání úloh, identifikátory studentů propojené s jejich řešeními nebo obsah diskuze v komentářích k odevzdání.
Způsob zpracování těchto dat, zejména při využití cloudových API, vyvolává závažné otázky z hlediska ochrany osobních údajů a souladu s GDPR\@.

\subsubsection*{GDPR a přenos dat do třetí strany}
\label{subsubsec:gdpr-data-transfer}

Použití cloudového API znamená, že vstupní data, včetně studentských zdrojových kódů a zadání úloh, jsou přenášena a zpracovávána na serverech externího poskytovatele, což vyvolává otázky souladu s nařízením GDPR\@.

Otázka, zda zdrojový kód studentů představuje osobní údaj ve smyslu čl.~4 odst.~1 GDPR\footnote{https://www.privacy-regulation.eu/cs/4.htm}, závisí vždy na konkrétním kontextu zpracování -- rozhodující není pouze obsah souboru, ale i způsob předání a dostupné prostředky identifikace.
Kód bez identifikátorů osobním údajem být nemusí, avšak v praxi zdrojové soubory často obsahují v záhlaví jméno autora, e-mail nebo login, případně jsou předávány spolu s metadaty.
V takovém případě jsou osobním údajem jednoznačně.
Toto bylo ověřeno prostřednictvím emailové komunikace s pověřencem pro ochranu osobních údajů VŠB-TUO\@.
Otázkou osobních údajů ve zdrojovém kódu se zabývá i odborná literatura~\cite{kocetkov-stack}.
Citlivější než samotný kód je pak bodové hodnocení, které je osobním údajem bezpochyby.

Pokud jsou před odesláním ze zdrojového kódu odstraněny všechny identifikátory a kód není předáván spolu s metadaty umožňujícími identifikaci studenta, ustanovení GDPR o přenosu osobních údajů do třetích zemí se neuplatní.
V opačném případě je nutné uzavřít smlouvu o zpracování osobních údajů (Data Processing Agreement, DPA)\footnote{https://gdpr.eu/what-is-data-processing-agreement/} a posoudit podmínky konkrétní služby, zejména zda poskytovatel data uchovává nebo využívá pro další trénování modelu.
Režimu GDPR lze tedy v~praxi předejít důslednou anonymizací vstupních dat před jejich odesláním do externího API\@.

\endinput